\documentclass[10pt,a4paper]{amsart}
\usepackage[margin=19mm]{geometry}
\usepackage{amsmath,amssymb,mathtools}
\usepackage[scr=rsfs,cal=euler]{mathalfa}
\usepackage{xfrac}
\usepackage[unicode,hidelinks,bookmarks=false]{hyperref}
\newcommand{\ie}{i.e.\ }
\newcommand{\cf}{cf.\ }
\newcommand{\resp}{respectively\ }
\newcommand{\Subsection}{\subsection}
\providecommand{\nobreakdash}{\nobreak\hbox{-}}
\setlength{\emergencystretch}{2em}
\multlinegap0pt
\usepackage{bm}
\usepackage{bbm}
\usepackage{dsfont}
\usepackage{xspace}
\usepackage{cancel}
\usepackage{mathtools}
\usepackage{amsthm}
\makeatletter
\@ifundefined{theorem}{\newtheorem{theorem}{Theorem}}{}
\@ifundefined{lemma}{\newtheorem{lemma}[theorem]{Lemma}}{}
\makeatother
\newcommand{\abs}[1]{\left|{#1}\right|}
\newcommand{\ds}{\mathrm{d}s}
\newcommand{\dtt}{\mathrm{d}_\tau}
\newcommand{\dt}{\mathrm{d}t}
\newcommand{\inpro}[2]{\left\langle#1,#2\right\rangle}
\newcommand{\norm}[2]{\left\|{#1}\right\|_{#2}}
\title[Error analysis of a fully discrete structure-preserving fini]{Error analysis of a fully discrete structure-preserving finite element scheme for a diffuse-interface model of tumour growth}
\author{Agus L. Soenjaya, Ping Lin, Thanh Tran}
\date{Source: arXiv:2509.14486v3 (CC BY 4.0)}
\begin{document}
\maketitle

\noindent\emph{Source and licence.} Error analysis of a fully discrete structure-preserving finite element scheme for a diffuse-interface model of tumour growth. Version arXiv:2509.14486v3 (CC BY 4.0), distributed under the Creative Commons Attribution 4.0 International licence, as recorded in the arXiv licence metadata for this exact version and shown on the abstract page https://arxiv.org/abs/2509.14486v3. Full rights notice: licenses/arxiv-2509.14486-CC-BY-4.0.txt.\par\smallskip

\noindent\textit{Editorial scope.} Counted unit: the Lemma whose source label is \texttt{equ:pre est N5} in the article (source0.tex lines 3253--3308, source label \texttt{equ:pre est N5}), delivered together with the article's own complete original proof of it (source0.tex lines 3310--3728, one \texttt{proof} environment of 419 lines). No mathematics is written, repaired or re-derived: statement and proof are the article's own text line for line. Required context retained uncounted: equ:fem euler r (lines 558--598); equ:phk bhk (lines 603--610); equ:beta k bk (lines 612--619); equ:stab u Linfty (lines 1684--1688); equ:stab dtu L2 (lines 1892--1907); equ:dt vn Lp (lines 2147--2153); equ:dt vn min par v (lines 2155--2162); equ:error euler L2 (lines 2520--2528). Every definition, display and label of those lines is retained unchanged. Omitted from this package and not redistributed: 1--557, 599--602, 611--611, 620--1683, 1689--1891, 1908--2146, 2154--2154, 2163--2519, 2529--3252, 3309--3309, 3729--5321. Nothing inside a retained excerpt is omitted or altered: every retained body is delivered exactly as the article's lines are written, so no omission receipt of any kind accompanies this package. Citations: the retained excerpts carry no citation command at all, so no bibliography entry of the article is transcribed and no reference is redistributed. Reference adaptation: none was needed and none was made; every reference inside the delivered unit resolves inside it, so no unresolved reference or citation is produced. Printed slips kept literal: no printed word, symbol, hypothesis or number of the counted statement or of its proof was altered. Layout and class: the article's own class is \texttt{amsart} with packages including amsmath, amsthm, amscd, amsfonts, amssymb, graphicx, color, mathtools, mathrsfs, hyperref, enumerate, setspace, multicol, geometry; the wrapper is the lane's pinned \texttt{amsart} editorial wrapper at 10pt on A4, which restates the theorem-like environments the retained text needs and adds the article's own macro definitions that the retained text needs (\texttt{\textbackslash abs}, \texttt{\textbackslash ds}, \texttt{\textbackslash dt}, \texttt{\textbackslash dtt}, \texttt{\textbackslash inpro}, \texttt{\textbackslash norm}), with extra wrapper packages (bm, bbm, dsfont, xspace, cancel, mathtools). The delivered numbering is the unit's own. Assets: no retained span contains a figure environment, a graphics-inclusion command, a PDF-page inclusion, a table or a TikZ picture; no image file is copied into this package, the package declares no asset dependency and no image file is redistributed with it. Licence: version arXiv:2509.14486v3 of this article is distributed under the Creative Commons Attribution 4.0 International licence, as recorded in the arXiv licence metadata for this exact version and shown on the abstract page https://arxiv.org/abs/2509.14486v3. A full rights notice is delivered at licenses/arxiv-2509.14486-CC-BY-4.0.txt. Characters: every retained excerpt is pure ASCII, so the delivered engine, XeLaTeX with its default Latin Modern text font, reports no missing character.
	\begin{alignat}{2}
		\label{equ:fem euler u}
		\inpro{\dtt u_h^{k+1}}{\phi}
		&=
		-\inpro{\nabla \mu_h^{k+1}}{\nabla \phi}
		+
		\inpro{P(u_h^k)\cdot (\sigma_h^{k+1}-\mu_h^{k+1})}{\phi}, \quad &&\forall \phi \in V_h,
		\\
		\inpro{\mu_h^{k+1}}{\psi}
		&=
		\epsilon^2 \inpro{\nabla u_h^{k+1}}{\nabla \psi}
		+
		\lambda \inpro{u_h^{k+1}}{\psi}
		-
		\chi_0 \inpro{n_h^{k+1}}{\psi}
		\nonumber\\
		\label{equ:fem euler mu}
		&\qquad
		+
		\frac{r_h^{k+1}}{\sqrt{\mathcal{E}_1[u_h^k]+B}} \inpro{f'(u_h^k)}{\psi}, \quad &&\forall \psi\in V_h,
		\\
		\label{equ:fem euler n}
		\inpro{\dtt n_h^{k+1}}{\varphi}
		&=
		-\inpro{\nabla \sigma_h^{k+1}}{\nabla \varphi}
		-
		\inpro{P(u_h^k)\cdot (\sigma_h^{k+1}-\mu_h^{k+1})}{\varphi}, \quad &&\forall \varphi \in V_h,
		\\
		\label{equ:fem euler sigma}
		\sigma_h^{k+1}
		&=
		\delta^{-1} n_h^{k+1}
		-
		\chi_0 u_h^k,
		\\
		\label{equ:fem euler r}
		\dtt r_h^{k+1}
		&=
		\frac{1}{2\sqrt{\mathcal{E}_1[u_h^k]+B}}
		\inpro{f'(u_h^k)}{\dtt u_h^{k+1}}.
	\end{alignat}

\begin{align}\label{equ:phk bhk}
	p_h^k:=P(u_h^k),
	\qquad
	\beta_h^k := \mathcal E_1[u_h^k]+B,
	\qquad
	b_h^k:=
	\frac{f'(u_h^k)}{\sqrt{\beta_h^k}},
\end{align}

\begin{align}\label{equ:beta k bk}
	p^k :=P(u^k),
	\qquad
	\beta^k := \mathcal E_1[u^k]+B = (r^k)^2,
	\qquad
	b^k:=
	\frac{f'(u^k)}{\sqrt{\beta^k}}.
\end{align}

\begin{align}\label{equ:stab u Linfty}
	\norm{u_h^k}{L^\infty}^2
    +
    \tau \sum_{j=1}^k \left(\norm{\mu_h^j}{L^\infty}^2 + \norm{n_h^j}{L^\infty}^2 + \norm{\sigma_h^j}{L^\infty}^2 \right) \leq C,
\end{align}

\begin{lemma}
For $k=1,2,\ldots, \lfloor T/\tau \rfloor$,
	\begin{align}\label{equ:stab dtu L2}
		\norm{\dtt u_h^k}{L^2}^2
		+
		\norm{\dtt n_h^k}{L^2}^2
		+
		\tau \sum_{j=1}^k \left(\norm{\nabla \dtt u_h^j}{L^2}^2 + \norm{\Delta_h \dtt u_h^j}{L^2}^2 + \norm{\dtt \mu_h^j}{L^2}^2 + \norm{\nabla \dtt n_h^j}{L^2}^2 \right) 
		\leq C.
	\end{align}
	In particular, this implies
	\begin{align}\label{equ:stab uk1 uk L2}
		\norm{u_h^{k+1}-u_h^k}{L^2} + \norm{n_h^{k+1}-n_h^k}{L^2} \leq C\tau,
	\end{align}
	where $C$ depends on $T$, but is independent of $k$, $h$, and $\tau$.
\end{lemma}

\begin{align}\label{equ:dt vn Lp}
	\norm{\dtt v^k}{L^p}
	&=
	\norm{\frac{1}{\tau} \int_{t_{k-1}}^{t_k} \partial_t v(t) \,\dt}{L^p}
	\leq
	\norm{v}{W^{1,\infty}_T(L^p)},
\end{align}

\begin{align}
	\label{equ:dt vn min par v}
	\norm{\dtt v^k- \partial_t v^k}{L^p}
	&=
	\norm{\frac{1}{2\tau} \int_{t_{k-1}}^{t_k} (t-t_{k-1}) \partial_t^2 v(t)\,\dt}{L^p}
	\leq 
	C_p \tau \norm{v}{W^{2,\infty}_T(L^p)}.
\end{align}

\begin{align}\label{equ:error euler L2}
	\norm{u_h^k-u(t_k)}{L^2}
	+
	\norm{n_h^k-n(t_k)}{L^2}
	+
	\abs{r_h^k-r(t_k)}
	&\leq
	C(h^{q+1}+\tau).
\end{align}

\begin{lemma}
	Let $k\in\{1,2,\ldots,\lfloor T/\tau\rfloor\}$. Let $\beta_h^k$, $b_h^k$, $\beta^k$, and $b^k$ be defined by \eqref{equ:phk bhk} and \eqref{equ:beta k bk}. For any $\alpha>0$, there exists a constant $C_\alpha>0$, independent of $h$, $\tau$, and $k$, such that for all $\psi\in V_h$,
	\begin{align}
		\label{equ:pre est N5}
		&\abs{
			\inpro{
				\frac{1}{\tau}
				\left(
				\frac{r_h^{k+1}}{\sqrt{\beta_h^k}}
				-
				\frac{r_h^k}{\sqrt{\beta_h^{k-1}}}
				\right)
				f'(u_h^k)
			}{\psi}}
		\nonumber\\
		&\qquad
		\leq
		C_\alpha\norm{\psi}{L^2}^2
		+
		\alpha h^{2(q+1)}
		+
		\alpha \tau^2
		+
		\alpha\left(
		\norm{\dtt\theta^{k+1}}{L^2}^2
		+
		\norm{\dtt\theta^k}{L^2}^2
		\right),
	\end{align}
	and
	\begin{align}
		\label{equ:pre est N6}
		&\abs{
			\inpro{
				\left(\frac{r_h^k}{\sqrt{\beta_h^{k-1}}}\right)
				\left(
				\frac{f'(u_h^k)-f'(u_h^{k-1})}{\tau}
				\right)
				-
				\left(
				\frac{r^{k+1}}{\sqrt{\beta^{k+1}}}
				\right)
				f''(u^{k+1})\partial_t u^{k+1}
			}{\psi}}
		\nonumber\\
		&\qquad
		\leq
		C_\alpha \norm{\psi}{L^2}^2
		+
		\alpha h^{2(q+1)}
		+
		\alpha \tau^2
		+
		\alpha\norm{\dtt\theta^k}{L^2}^2.
	\end{align}
\end{lemma}

\begin{proof}
	We first prove \eqref{equ:pre est N5}. By \eqref{equ:fem euler r}, we have
	\begin{align}
		\frac{1}{\tau}
		\left(
		\frac{r_h^{k+1}}{\sqrt{\beta_h^k}}
		-
		\frac{r_h^k}{\sqrt{\beta_h^{k-1}}}
		\right)
		&=
		\frac{\dtt r_h^{k+1}}{\sqrt{\beta_h^k}}
		+
		\frac{1}{\tau}
		\left(
		\frac{1}{\sqrt{\beta_h^k}}
		-
		\frac{1}{\sqrt{\beta_h^{k-1}}}
		\right) r_h^k
		\nonumber\\
		&\quad
		=
		\frac{1}{2\beta_h^k}
		\inpro{f'(u_h^k)}{\dtt u_h^{k+1}}
		-
		\frac{r_h^k}{\sqrt{\beta_h^k\beta_h^{k-1}}
			(\sqrt{\beta_h^k}+\sqrt{\beta_h^{k-1}})}
		\dtt\beta_h^k .
		\label{equ:ratio-diff-rh}
	\end{align}
	Using the integral form of the mean-value theorem, we write
	\[
	\dtt\beta_h^k
	=
	\inpro{\overline f_h^k}{\dtt u_h^k},
	\qquad \text{where}\quad
	\overline f_h^k
	:=
	\int_0^1
	f'\big(u_h^{k-1}+s(u_h^k-u_h^{k-1})\big)\,\ds .
	\]
	Hence \eqref{equ:ratio-diff-rh} becomes
	\begin{align}
		&\frac{1}{\tau}
		\left(
		\frac{r_h^{k+1}}{\sqrt{\beta_h^k}}
		-
		\frac{r_h^k}{\sqrt{\beta_h^{k-1}}}
		\right)
		=
		\frac{1}{2\beta_h^k}
		\inpro{f'(u_h^k)}{\dtt u_h^{k+1}}
		-
		\frac{r_h^k}{\sqrt{\beta_h^k\beta_h^{k-1}}
			(\sqrt{\beta_h^k}+\sqrt{\beta_h^{k-1}})}
		\inpro{\overline f_h^k}{\dtt u_h^k}.
		\label{equ:ratio-diff-rh-2}
	\end{align}
	We now estimate the two terms on the right-hand side of
	\eqref{equ:ratio-diff-rh-2} separately. Adding and subtracting a common exact
	quantity
	\[
	\frac{1}{2\beta^k}
	\inpro{f'(u^k)}{\partial_t u^k},
	\]
	we obtain by the triangle inequality,
	\begin{align}
		&\abs{
			\frac{1}{\tau}
			\left(
			\frac{r_h^{k+1}}{\sqrt{\beta_h^k}}
			-
			\frac{r_h^k}{\sqrt{\beta_h^{k-1}}}
			\right)}
		\nonumber\\
		&\quad
		\leq
		\abs{
			\frac{1}{2\beta_h^k}
			\inpro{f'(u_h^k)}{\dtt u_h^{k+1}}
			-
			\frac{1}{2\beta^k}
			\inpro{f'(u^k)}{\partial_t u^k}
		}
		\nonumber\\
		&\qquad
		+
		\abs{
			\frac{r_h^k}{\sqrt{\beta_h^k\beta_h^{k-1}}
				(\sqrt{\beta_h^k}+\sqrt{\beta_h^{k-1}})}
			\inpro{\overline f_h^k}{\dtt u_h^k}
			-
			\frac{1}{2\beta^k}
			\inpro{f'(u^k)}{\partial_t u^k}
		}
		\nonumber\\
		&\quad
		=: J_1+J_2 .
		\label{equ:J1-J2-ratio}
	\end{align}
	For the first term, we have by adding and subtracting common terms,
	\begin{align}
		J_1
		&\leq
		\abs{
			\frac{1}{2\beta_h^k}
			-
			\frac{1}{2\beta^k}
		}
		\abs{\inpro{f'(u_h^k)}{\dtt u_h^{k+1}}}
		\nonumber\\
		&\quad
		+
		C\abs{
			\inpro{f'(u_h^k)-f'(u^k)}{\dtt u_h^{k+1}}
		}
		+
		C\abs{
			\inpro{f'(u^k)}{\dtt u_h^{k+1}-\partial_t u^k}
		}.
		\label{equ:J1-est-decomp}
	\end{align}
	Now, observe that we have
	\[
	\abs{\beta_h^k-\beta^k} + \norm{f'(u_h^k)-f'(u^k)}{L^2}
	\leq
	C\norm{u_h^k-u^k}{L^2}
	\leq
	C(h^{q+1}+\tau),
	\]
	as well as \eqref{equ:stab u Linfty}, \eqref{equ:stab dtu L2}, \eqref{equ:dt vn Lp}, and \eqref{equ:dt vn min par v}. Thus, we deduce from \eqref{equ:J1-est-decomp} that
	\begin{align}
		J_1
		&\leq
		C(h^{q+1}+\tau)
		+
		C\norm{\dtt u_h^{k+1}-\partial_t u^k}{L^2}
		\nonumber\\
		&\leq
		C(h^{q+1}+\tau)
		+
		C\norm{\dtt\theta^{k+1}+
			\dtt\rho^{k+1}
			+
			\dtt u^{k+1}-\partial_t u^k}{L^2}
		\nonumber\\
		&\leq
		C(h^{q+1}+\tau)
		+
		C\norm{\dtt\theta^{k+1}}{L^2}.
		\label{equ:J1-est}
	\end{align}
	We next estimate $J_2$. First, we note the following fact: the map
	\[
	G(r,a,b):=\frac{r}{\sqrt{ab}(\sqrt a+\sqrt b)}
	\]
	is Lipschitz on the range where $a,b$ are bounded away from zero and $r$ is bounded. We also observe that
	$G(r^k,\beta^k,\beta^k)=1/(2\beta^k)$. Therefore,
	\begin{align}
		&\abs{
			\frac{r_h^k}{\sqrt{\beta_h^k\beta_h^{k-1}}
				(\sqrt{\beta_h^k}+\sqrt{\beta_h^{k-1}})}
			-
			\frac{1}{2\beta^k}
		}
		\nonumber\\
		&=
		\abs{G(r_h^k,\beta_h^k,\beta_h^{k-1})- G(r^k,\beta^k,\beta^k)}
		\nonumber\\
		&
		\leq
		C\abs{r_h^k-r^k}
		+
		C\abs{\beta_h^k-\beta^k}
		+
		C\abs{\beta_h^{k-1}-\beta^{k-1}}
		+
		C\abs{\beta^k-\beta^{k-1}}
		\nonumber\\
		&
		\leq
		C(h^{q+1}+\tau)
		\label{equ:coef-J2}
	\end{align}
	by \eqref{equ:error euler L2}.
	Moreover, by the definition of $\overline f_h^k$ and the mean-value theorem,
	\begin{align}
		\norm{\overline f_h^k-f'(u^k)}{L^2}
		&\leq
		C\norm{u_h^k-u^k}{L^2}
		+
		C\norm{u_h^{k-1}-u^k}{L^2}
		\leq
		C(h^{q+1}+\tau).
		\label{equ:fbar-est}
	\end{align}
	Thus, we have the estimate
	\begin{align}
		J_2
		&\leq
		C(h^{q+1}+\tau)
		\abs{\inpro{\overline f_h^k}{\dtt u_h^k}}
		+
		C\abs{\inpro{\overline f_h^k-f'(u^k)}{\dtt u_h^k}}
		+
		C\abs{\inpro{f'(u^k)}{\dtt u_h^k-\partial_t u^k}}
		\nonumber\\
		&\leq
		C(h^{q+1}+\tau)
		+
		C\norm{\dtt u_h^k-\partial_t u^k}{L^2}
		\nonumber\\
		&\leq
		C(h^{q+1}+\tau)
		+
		C\norm{\dtt\theta^k}{L^2}.
		\label{equ:J2-est}
	\end{align}
	Combining \eqref{equ:J1-J2-ratio}, \eqref{equ:J1-est}, and
	\eqref{equ:J2-est}, we obtain
	\begin{align}
		\label{equ:ratio-pre-est}
		\abs{
			\frac{1}{\tau}
			\left(
			\frac{r_h^{k+1}}{\sqrt{\beta_h^k}}
			-
			\frac{r_h^k}{\sqrt{\beta_h^{k-1}}}
			\right)}
		\leq
		C(h^{q+1}+\tau)
		+
		C\norm{\dtt\theta^{k+1}}{L^2}
		+
		C\norm{\dtt\theta^k}{L^2}.
	\end{align}
	By \eqref{equ:stab u Linfty}, we also have
	$\norm{f'(u_h^k)}{L^\infty}\le C$. Inequality \eqref{equ:pre est N5} then follows by an application of Young's inequality.

	We now prove \eqref{equ:pre est N6}. Again by the integral form of the mean-value theorem,
	\[
	\frac{f'(u_h^k)-f'(u_h^{k-1})}{\tau}
	=
	\overline f_{h,2}^k\, \dtt u_h^k, \qquad \text{where} \quad 
	\overline f_{h,2}^k
	:=
	\int_0^1
	f''\big(u_h^{k-1}+s(u_h^k-u_h^{k-1})\big)\,\ds .
	\]
	We first record the following estimate.
	Adding and subtracting
	$\overline f_{h,2}^k\partial_tu^{k+1}$, we have
	\begin{align*}
		&\norm{
			\frac{f'(u_h^k)-f'(u_h^{k-1})}{\tau}
			-
			f''(u^{k+1})\partial_tu^{k+1}
		}{L^2}
		\nonumber\\
		&=
		\norm{
			\overline f_{h,2}^k \dtt u_h^k
			-
			f''(u^{k+1})\partial_tu^{k+1}
		}{L^2}
		\nonumber\\
		&
		\leq
		\norm{\overline f_{h,2}^k}{L^\infty}
		\norm{\dtt u_h^k-\partial_tu^{k+1}}{L^2}
		+
		\norm{\overline f_{h,2}^k-f''(u^{k+1})}{L^2}
		\norm{\partial_tu^{k+1}}{L^\infty}
		\nonumber\\
		&=: L_1+L_2.
	\end{align*}
	By the uniform $L^\infty$ bound for $u_h^j$ in \eqref{equ:stab u Linfty}, the growth assumption on $f''$, \eqref{equ:dt vn Lp}, and \eqref{equ:dt vn min par v}, we have
	\begin{align}
		L_1
		&\leq
		C\left(\norm{\dtt\theta^k}{L^2}
		+
		\norm{\dtt\rho^k}{L^2}
		+
		\norm{\dtt u^k-\partial_tu^{k+1}}{L^2}\right)
		\leq
		C\norm{\dtt\theta^k}{L^2}
		+
		C(h^{q+1}+\tau).
		\label{equ:N6-dtuh-est}
	\end{align}
	Moreover, by the mean-value theorem, the assumed regularity on $u$, and \eqref{equ:error euler L2}, we have
	\begin{align}
		L_2
		&\leq
		C\norm{u_h^k-u^{k+1}}{L^2}
		+
		C\norm{u_h^{k-1}-u^{k+1}}{L^2}
		\leq
		C(h^{q+1}+\tau).
		\label{equ:N6-fbar2-est}
	\end{align}
	Combining \eqref{equ:N6-dtuh-est}--\eqref{equ:N6-fbar2-est}, we obtain
	\begin{align}
		\label{equ:diff-fprime-est}
		&\norm{
			\frac{f'(u_h^k)-f'(u_h^{k-1})}{\tau}
			-
			f''(u^{k+1})\partial_tu^{k+1}
		}{L^2}
		\leq
		C(h^{q+1}+\tau)
		+
		C\norm{\dtt\theta^k}{L^2}.
	\end{align}
	
	Next, since $r^{k+1}/\sqrt{\beta^{k+1}}=1$ by definition, we have the estimate
	\begin{align}
		\abs{
			\frac{r_h^k}{\sqrt{\beta_h^{k-1}}}
			-
			\frac{r^{k+1}}{\sqrt{\beta^{k+1}}}
		}
		&
		=
		\abs{
			\frac{r_h^k}{\sqrt{\beta_h^{k-1}}}
			-
			\frac{r^k}{\sqrt{\beta^k}}
		}
		\nonumber\\
		&\quad
		\leq
		C\abs{r_h^k-r^k}
		+
		C\abs{\beta_h^{k-1}-\beta^k}
		\nonumber\\
		&\quad
		\leq
		C\abs{r_h^k-r^k}
		+
		C\abs{\beta_h^{k-1}-\beta^{k-1}}
		+
		C\abs{\beta^{k-1}-\beta^k}
		\nonumber\\
		&\quad
		\leq
		C(h^{q+1}+\tau).
		\label{equ:N6-ratio-est}
	\end{align}
	Therefore, by adding and subtracting a common term
	\[
	\frac{r^{k+1}}{\sqrt{\beta^{k+1}}}
	\frac{f'(u_h^k)-f'(u_h^{k-1})}{\tau},
	\]
	we obtain
	\begin{align}
		&\norm{
			\left(\frac{r_h^k}{\sqrt{\beta_h^{k-1}}}\right)
			\left(
			\frac{f'(u_h^k)-f'(u_h^{k-1})}{\tau}
			\right)
			-
			\left(
			\frac{r^{k+1}}{\sqrt{\beta^{k+1}}}
			\right)
			f''(u^{k+1})\partial_tu^{k+1}
		}{L^2}
		\nonumber\\
		&\quad
		\leq
		\abs{
			\frac{r_h^k}{\sqrt{\beta_h^{k-1}}}
			-
			\frac{r^{k+1}}{\sqrt{\beta^{k+1}}}
		}
		\norm{
			\frac{f'(u_h^k)-f'(u_h^{k-1})}{\tau}
		}{L^2}
		\nonumber\\
		&\qquad
		+
		C
		\norm{
			\frac{f'(u_h^k)-f'(u_h^{k-1})}{\tau}
			-
			f''(u^{k+1})\partial_tu^{k+1}
		}{L^2}.
		\label{equ:N6-product-split}
	\end{align}
	Now, we note that by \eqref{equ:stab dtu L2}, the growth assumption on $f''$, and \eqref{equ:stab u Linfty},
	\[
	\norm{
		\frac{f'(u_h^k)-f'(u_h^{k-1})}{\tau}
	}{L^2}
	=
	\norm{\overline f_{h,2}^k\, \dtt u_h^k}{L^2}
	\leq 
	C,
	\]
	Therefore, using \eqref{equ:diff-fprime-est} and \eqref{equ:N6-ratio-est}, we infer from \eqref{equ:N6-product-split} that
	\begin{align*}
		&\norm{
			\left(\frac{r_h^k}{\sqrt{\beta_h^{k-1}}}\right)
			\left(
			\frac{f'(u_h^k)-f'(u_h^{k-1})}{\tau}
			\right)
			-
			\left(
			\frac{r^{k+1}}{\sqrt{\beta^{k+1}}}
			\right)
			f''(u^{k+1})\partial_tu^{k+1}
		}{L^2}
		\leq
		C(h^{q+1}+\tau)
		+
		C\norm{\dtt\theta^k}{L^2}.
	\end{align*}
	We then deduce \eqref{equ:pre est N6} by an application of Young's inequality. The proof is now complete.
\end{proof}

\end{document}
